\documentclass[a4paper,12pt]{ctexart}
\usepackage{xcolor}
\usepackage[top=1.8cm, bottom=1.8cm, left=1.8cm, right=1.8cm]{geometry}
\usepackage{array}
\usepackage{listings}
\usepackage{hyperref}
\usepackage{amsmath, amssymb}
\hypersetup{colorlinks=true,linkcolor=blue!70!black}
\linespread{1.35}

\lstset{
  basicstyle=\ttfamily\small,
  keywordstyle=\color{blue},
  commentstyle=\color{gray},
  showstringspaces=false,
  frame=single,
  breaklines=true
}

\title{\textcolor{blue!70!black}{\Huge\textbf{计算机数学基础}}}
\author{编程背后的数学}
\date{}

\begin{document}
\maketitle
\thispagestyle{empty}

\section*{写在前面}

数学和编程是密不可分的。你可能在学校里学过一些数学，但没想过它们和编程有什么关系。这个模块会把\textbf{对编程最有用的数学知识}挑出来，用代码来理解数学，用数学来优化代码。

\begin{center}
\begin{tabular}{|c|c|}
\hline
\textbf{本章内容} & \textbf{用在编程的哪里} \\ \hline
二进制与位运算 & 数据存储、权限系统、性能优化 \\ \hline
布尔代数 & 条件判断、搜索过滤 \\ \hline
集合论 & 数据处理、去重、数据库查询 \\ \hline
模运算 & 哈希函数、加密、循环调度 \\ \hline
组合数学 & 算法复杂度分析、概率算法 \\ \hline
图论基础 & 地图导航、社交网络、推荐系统 \\ \hline
逻辑与证明 & 程序正确性、调试思维 \\ \hline
\end{tabular}
\end{center}

\section*{第一课：二进制——计算机的母语}

\subsection*{为什么是二进制？}

计算机内部有数十亿个微小开关（晶体管），每个开关只有两种状态：\textbf{开（1）}或\textbf{关（0）}。所以计算机只懂二进制。

\subsection*{十进制 vs 二进制}

\begin{center}
\begin{tabular}{|c|c|c|}
\hline
\textbf{十进制} & \textbf{二进制} & \textbf{解释} \\ \hline
0 & 0 & $0 \times 2^0 = 0$ \\ \hline
1 & 1 & $1 \times 2^0 = 1$ \\ \hline
2 & 10 & $1 \times 2^1 + 0 \times 2^0 = 2$ \\ \hline
3 & 11 & $1 \times 2^1 + 1 \times 2^0 = 3$ \\ \hline
4 & 100 & $1 \times 2^2 + 0 + 0 = 4$ \\ \hline
5 & 101 & $1 \times 2^2 + 0 + 1 = 5$ \\ \hline
8 & 1000 & $1 \times 2^3 = 8$ \\ \hline
15 & 1111 & $1+2+4+8 = 15$ \\ \hline
\end{tabular}
\end{center}

\textbf{规律：} 从右往左，每一位代表 $2^0, 2^1, 2^2, 2^3 \dots$

\subsection*{二进制转十进制}

\textbf{方法：} 把每一位的值（0或1）乘以 $2^{\text{位置}}$，然后加起来。

\begin{lstlisting}[language=Python]
def binary_to_decimal(bin_str):
    """把二进制字符串转十进制，如 '1011' → 11"""
    result = 0
    for i, digit in enumerate(reversed(bin_str)):
        if digit == '1':
            result += 2 ** i
    return result

print(binary_to_decimal("1011"))  # 11
print(binary_to_decimal("1111"))  # 15
\end{lstlisting}

\subsection*{十进制转二进制}

\textbf{方法：} 不断除以 2，余数倒着读。

\begin{lstlisting}[language=Python]
def decimal_to_binary(n):
    """把十进制转二进制字符串，如 11 → '1011'"""
    if n == 0:
        return "0"
    result = ""
    while n > 0:
        result = str(n % 2) + result
        n = n // 2
    return result

for n in [0, 1, 2, 5, 10, 15, 42]:
    print(f"{n:3d} → {decimal_to_binary(n)}")
\end{lstlisting}

\subsection*{位运算——让程序飞起来}

位运算直接操作二进制位，比普通数学运算快得多。

\begin{lstlisting}[language=Python]
a = 5   # 二进制 0101
b = 3   # 二进制 0011

print(a & b)   # 按位与 AND → 0001 = 1
print(a | b)   # 按位或 OR  → 0111 = 7
print(a ^ b)   # 按位异或 XOR → 0110 = 6
print(~a)      # 按位取反 NOT → -(5+1) = -6
print(a << 1)  # 左移1位 → 1010 = 10（相当于 ×2）
print(a >> 1)  # 右移1位 → 0010 = 2（相当于 ÷2）
\end{lstlisting}

\textbf{实用技巧：}
\begin{itemize}
  \item \texttt{n \& 1} 判断奇偶（比 \texttt{n \% 2} 快）：结果为 1 是奇数，0 是偶数
  \item \texttt{n << 1} 相当于 \texttt{n * 2}
  \item \texttt{n >> 1} 相当于 \texttt{n // 2}
  \item 用位运算实现权限系统：每个 bit 代表一个权限
\end{itemize}

\subsection*{练习}

\begin{enumerate}
  \item 把二进制 \texttt{11010} 转十进制，把十进制 \texttt{27} 转二进制
  \item 用 Python 验证：\texttt{5 \& 3} 的结果和手动计算是否一致
  \item 用位运算判断一个数是否是 2 的幂（提示：2 的幂的二进制只有一个 1）
  \item \textbf{[扩展]} 实现一个简单的权限系统：读取=1，写入=2，执行=4，用 \texttt{|} 组合权限，用 \texttt{\&} 检查权限
\end{enumerate}

\section*{第二课：八进制与十六进制}

二进制写起来太长了（比如 255 是 11111111），所以程序员常用\textbf{八进制}（每 3 位二进制一组）和\textbf{十六进制}（每 4 位一组）。

\begin{center}
\begin{tabular}{|c|c|c|c|}
\hline
\textbf{十进制} & \textbf{二进制} & \textbf{八进制} & \textbf{十六进制} \\ \hline
0 & 0000 & 0 & 0 \\ \hline
1 & 0001 & 1 & 1 \\ \hline
10 & 1010 & 12 & A \\ \hline
15 & 1111 & 17 & F \\ \hline
16 & 10000 & 20 & 10 \\ \hline
255 & 11111111 & 377 & FF \\ \hline
\end{tabular}
\end{center}

\textbf{十六进制在编程中到处可见：}
\begin{itemize}
  \item 颜色：\texttt{\#FF5733}（红色=FF, 绿色=57, 蓝色=33）
  \item 内存地址：\texttt{0x7ffd5e3a}
  \item Unicode 字符：\texttt{U+4E2D}（"中"字）
\end{itemize}

\begin{lstlisting}[language=Python]
# Python 中的进制表示
print(0b1010)     # 二进制 → 10
print(0o12)       # 八进制 → 10
print(0xA)        # 十六进制 → 10

# 进制转换
print(bin(42))     # '0b101010'
print(oct(42))     # '0o52'
print(hex(42))     # '0x2a'
print(int("FF", 16))  # 255
\end{lstlisting}

\subsection*{练习}

\begin{enumerate}
  \item 把十进制 255 转成八进制和十六进制。用 Python 验证。
  \item 颜色 \texttt{\#AABBCC} 的红色、绿色、蓝色分量分别是多少（转十进制）？
  \item \textbf{[扩展]} 查一下为什么程序员常用十六进制而不是八进制来表示颜色和内存地址？
\end{enumerate}

\section*{第三课：日常生活中的数学概念映射}

当你学编程时，很多概念其实就是你熟悉的数学概念的"另一种说法"：

\begin{center}
\begin{tabular}{|c|c|}
\hline
\textbf{数学} & \textbf{编程} \\ \hline
函数 $f(x) = x + 1$ & \texttt{def f(x): return x + 1} \\ \hline
集合 $\{1, 2, 3\}$ & \texttt{\{1, 2, 3\}} \\ \hline
不等式 $x > 0$ & \texttt{x > 0} \\ \hline
求和 $\sum_{i=1}^{n} i$ & \texttt{sum(range(1, n+1))} \\ \hline
数列 $a_1, a_2, \dots, a_n$ & \texttt{lst[0], lst[1], ..., lst[n-1]} \\ \hline
坐标 $(x, y)$ & \texttt{(x, y)} 元组 \\ \hline
逻辑与/或/非 ($\land, \lor, \lnot$) & \texttt{and, or, not} \\ \hline
\end{tabular}
\end{center}

\subsection*{练习}

\begin{enumerate}
  \item 写出以下数学表达式的 Python 代码：① $f(x) = 2x + 3$ ② 集合 $\{1, 2, 3\}$ 与 $\{3, 4, 5\}$ 的交集 ③ $\sum_{i=1}^{10} i^2$
  \item Python 中坐标为 $(3, 7)$ 的点，数学上可以怎么表示？Python 可以用什么数据类型表示？
  \item \textbf{[思考]} 为什么 Python 用 \texttt{and} / \texttt{or} / \texttt{not} 而不是 $\land$ / $\lor$ / $\lnot$？这体现了编程语言的什么设计原则？
\end{enumerate}

\section*{第四课：模运算——时钟算术}

\subsection*{什么是模运算？}

模运算就是\textbf{求余数}。生活中处处是模运算：现在是 10 点，过 5 小时是几点？$10 + 5 = 15 \equiv 3 \pmod{12}$。

\begin{lstlisting}[language=Python]
# Python 中的模运算
print(17 % 5)    # 2 (17÷5=3 余 2)
print(20 % 7)    # 6 (20÷7=2 余 6)
print(-3 % 7)    # 4 (Python 特殊：结果总为正)
\end{lstlisting}

\subsection*{模运算的应用}

\textbf{1. 判断整除}

\begin{lstlisting}[language=Python]
def is_even(n):   return n % 2 == 0
def is_odd(n):    return n % 2 == 1
def is_leap_year(y): return (y % 4 == 0 and y % 100 != 0) or (y % 400 == 0)

print(is_leap_year(2024))  # True
print(is_leap_year(2100))  # False
\end{lstlisting}

\textbf{2. 循环调度}

\begin{lstlisting}[language=Python]
# 轮流分配给 3 个玩家
players = ["小明", "小红", "小刚"]
for round in range(10):
    player = players[round % 3]
    print(f"第{round+1}轮：{player}的回合")
\end{lstlisting}

\textbf{3. 哈希函数}

\begin{lstlisting}[language=Python]
# 简单的哈希：把字符串映射到 0-99
def simple_hash(text):
    total = 0
    for ch in text:
        total += ord(ch)
    return total % 100

print(simple_hash("hello"))    # 某个 0-99 的数
print(simple_hash("world"))    # 不同的数
\end{lstlisting}

\textbf{4. 大数计算}（计算机竞赛常用）

\begin{lstlisting}[language=Python]
# 计算 2^1000 的最后三位
print(pow(2, 1000, 1000))  # 376（快速幂 + 模运算）
# 比先算 2^1000 再 % 1000 快无数倍
\end{lstlisting}

\subsection*{练习}

\begin{enumerate}
  \item 今天是星期三，100 天后是星期几？
  \item 写一个函数判断一个年份是否是闰年
  \item \textbf{[扩展]} 实现循环队列：用模运算实现固定大小数组的循环读写
\end{enumerate}

\section*{第五课：布尔代数——逻辑运算}

\subsection*{基本逻辑门}

布尔代数只有三个基本运算，但可以组合出任何逻辑：

\begin{center}
\begin{tabular}{|c|c|c|c|c|}
\hline
\textbf{A} & \textbf{B} & \textbf{A AND B} & \textbf{A OR B} & \textbf{NOT A} \\ \hline
0 & 0 & 0 & 0 & 1 \\ \hline
0 & 1 & 0 & 1 & 1 \\ \hline
1 & 0 & 0 & 1 & 0 \\ \hline
1 & 1 & 1 & 1 & 0 \\ \hline
\end{tabular}
\end{center}

\subsection*{用 Python 验证}

\begin{lstlisting}[language=Python]
# 真值表验证
for a in [True, False]:
    for b in [True, False]:
        print(f"{a:5} AND {b:5} = {a and b:5}")
        
# XOR（异或）：两个相同为假，不同为真
def xor(a, b):
    return (a and not b) or (not a and b)

print(xor(True, False))   # True
print(xor(True, True))    # False
\end{lstlisting}

\subsection*{德摩根定律——简化条件}

\[
\lnot (A \land B) = \lnot A \lor \lnot B
\]
\[
\lnot (A \lor B) = \lnot A \land \lnot B
\]

\textbf{编程中的应用：}

\begin{lstlisting}[language=Python]
# 原始条件：不在 x>0 且 y>0 范围内
if not (x > 0 and y > 0):
    print("至少有一个不大于0")

# 用德摩根简化后
if x <= 0 or y <= 0:
    print("至少有一个不大于0")
\end{lstlisting}

\subsection*{练习}

\begin{enumerate}
  \item 写出 XOR 的真值表
  \item 用德摩根定律简化：\texttt{not (score >= 60 or grade == "F")}
  \item \textbf{[扩展]} 只用 AND 和 NOT 实现 OR 的功能
\end{enumerate}

\section*{第六课：集合论——并集、交集、差集}

\subsection*{集合在编程中的对应}

数学中的集合论直接对应 Python 的 \texttt{set} 类型。

\begin{lstlisting}[language=Python]
A = {1, 2, 3, 4, 5}
B = {4, 5, 6, 7, 8}

print(A | B)   # 并集：{1,2,3,4,5,6,7,8}
print(A & B)   # 交集：{4, 5}
print(A - B)   # 差集（A有B没有）：{1, 2, 3}
print(A ^ B)   # 对称差（不同时有）：{1,2,3,6,7,8}
\end{lstlisting}

\subsection*{编程中的实际应用}

\begin{lstlisting}[language=Python]
# 场景：找出同时选了 Python 和 C++ 课的学生
python_students = {"小明", "小红", "小刚", "小李"}
cpp_students = {"小红", "小刚", "小王", "小赵"}

both = python_students & cpp_students
only_python = python_students - cpp_students
all_students = python_students | cpp_students

print(f"都选的：{both}")
print(f"只选Python的：{only_python}")
print(f"总共学生：{len(all_students)}人")

# 子集判断
print({1, 2}.issubset({1, 2, 3}))   # True
print({1, 2, 3}.issuperset({1, 2})) # True
\end{lstlisting}

\subsection*{练习}

\begin{enumerate}
  \item 有 3 个列表：\texttt{A = [1,2,3,4,5]}, \texttt{B = [4,5,6,7]}, \texttt{C = [1,3,5,7]}，找出只出现在一个列表中的元素
  \item 用集合判断两个字符串是否有相同字母
  \item \textbf{[扩展]} 用集合实现搜索引擎的"与或非"查询
\end{enumerate}

\section*{第七课：组合数学——数一数有多少种可能}

\subsection*{加法原理和乘法原理}

\begin{itemize}
  \item \textbf{加法原理：} 做一件事有 3 种方法，做另一件事有 4 种方法，总共 3+4=7 种
  \item \textbf{乘法原理：} 第一件事有 3 种选法，第二件有 4 种，总共 3×4=12 种
\end{itemize}

\subsection*{排列与组合}

\begin{lstlisting}[language=Python]
import math
from itertools import permutations, combinations

# 排列：从 3 个元素中选 2 个，考虑顺序
items = ['A', 'B', 'C']
print(list(permutations(items, 2)))
# [('A','B'), ('A','C'), ('B','A'), ('B','C'), ('C','A'), ('C','B')]
# 公式：P(3,2) = 3!/1! = 6

# 组合：从 3 个中选 2 个，不考虑顺序
print(list(combinations(items, 2)))
# [('A','B'), ('A','C'), ('B','C')]
# 公式：C(3,2) = 3!/(2!×1!) = 3

# Python 数学库直接算
print(math.perm(5, 3))     # 60
print(math.comb(5, 3))     # 10
\end{lstlisting}

\subsection*{实用场景}

\textbf{密码强度评估：}

\begin{lstlisting}[language=Python]
import math

# 6 位纯数字密码
digits = 10 ** 6
print(f"6位数字密码：{digits:,} 种组合")

# 8 位混合密码（大小写+数字+符号=94个字符）
chars = 94 ** 8
print(f"8位混合密码：{chars:,} 种组合")

# 从 50 人中选 3 人做班委
print(f"选班委方案：{math.comb(50, 3):,} 种")
\end{lstlisting}

\subsection*{练习}

\begin{enumerate}
  \item 4 位数字 PIN 码有多少种可能？
  \item 从 10 本不同的书中选 3 本，有多少种选法？
  \item \textbf{[扩展]} 写一个程序生成 ${}_6\text{C}_3$ 的所有组合，输出到文件
\end{enumerate}

\section*{第八课：图论基础——连接的世界}

\subsection*{图的基本概念}

图由\textbf{顶点}（Vertex）和\textbf{边}（Edge）组成。

\begin{lstlisting}[language=Python]
# 用邻接表表示图
graph = {
    "北京": ["天津", "石家庄", "上海"],
    "天津": ["北京", "唐山"],
    "石家庄": ["北京", "太原"],
    "上海": ["北京", "杭州", "南京"],
    "太原": ["石家庄"],
    "唐山": ["天津"],
    "杭州": ["上海"],
    "南京": ["上海"]
}

# 顶点数
print(f"顶点数：{len(graph)}")

# 边数
edges = sum(len(neighbors) for neighbors in graph.values()) // 2
print(f"边数：{edges}")

# 每个顶点的度（连接数）
for city, neighbors in graph.items():
    print(f"{city}：{len(neighbors)} 条连接")
\end{lstlisting}

\subsection*{树——没有环的图}

树是图的一种特殊情况：没有环，且所有顶点都连通。

\textbf{文件系统就是一棵树：}

\begin{lstlisting}[language=Python]
# 用嵌套字典表示树
file_system = {
    "name": "/",
    "type": "dir",
    "children": [
        {"name": "home", "type": "dir", "children": [
            {"name": "user", "type": "dir", "children": [
                {"name": "documents", "type": "dir", "children": [
                    {"name": "readme.md", "type": "file", "size": 1024}
                ]}
            ]}
        ]},
        {"name": "usr", "type": "dir", "children": [
            {"name": "bin", "type": "dir", "children": []}
        ]}
    ]
}

def count_files(node):
    if node["type"] == "file":
        return 1
    return sum(count_files(child) for child in node["children"])

print(f"文件数：{count_files(file_system)}")
\end{lstlisting}

\subsection*{图论的经典问题}

\textbf{最短路径：} 用 Dijkstra 算法（在 ds-algo 模块中有实现）

\textbf{最小生成树：} 用 Kruskal 或 Prim 算法

\textbf{拓扑排序：} 先修课程安排

\subsection*{练习}

\begin{enumerate}
  \item 画一个 5 个顶点的完全图（每两个顶点之间都有边），有几条边？
  \item 用邻接表表示一个社交网络（5个人，好友关系自己定义）
  \item \textbf{[扩展]} 计算任意两个城市之间的最短路径（用 BFS 或 Dijkstra）
\end{enumerate}

\section*{第九课：数列与求和}

\subsection*{等差数列}

\begin{lstlisting}[language=Python]
# 1 + 2 + 3 + ... + 100
# 公式：S = n * (a1 + an) / 2
def arithmetic_sum(a1, an, n):
    return n * (a1 + an) // 2

print(arithmetic_sum(1, 100, 100))  # 5050

# 用循环验证
print(sum(range(1, 101)))  # 5050
\end{lstlisting}

\subsection*{等比数列}

\begin{lstlisting}[language=Python]
# 1 + 2 + 4 + 8 + ... + 2^n
# 公式：a1 * (q^n - 1) / (q - 1)
def geometric_sum(a1, q, n):
    return a1 * (q ** n - 1) // (q - 1)

print(geometric_sum(1, 2, 10))  # 1023
# 验证：1+2+4+8+16+32+64+128+256+512 = 1023
\end{lstlisting}

\subsection*{练习}

\begin{enumerate}
  \item 计算 1 到 1000 所有奇数的和
  \item 二分查找最多能查多少次？$2^k \ge n$，解出 $k$——这和什么数列有关？
  \item \textbf{[扩展]} 等比数列求和公式推导：为什么 $\frac{a_1(q^n-1)}{q-1}$？
\end{enumerate}

\section*{第十课：矩阵——表格化的数据}

矩阵在 AI/机器学习、图形学、游戏开发中无处不在。

\begin{lstlisting}[language=Python]
# 用嵌套列表表示矩阵
def create_matrix(rows, cols, default=0):
    return [[default] * cols for _ in range(rows)]

def matrix_add(A, B):
    """矩阵加法"""
    rows, cols = len(A), len(A[0])
    C = create_matrix(rows, cols)
    for i in range(rows):
        for j in range(cols):
            C[i][j] = A[i][j] + B[i][j]
    return C

def matrix_multiply(A, B):
    """矩阵乘法"""
    rows, cols = len(A), len(B[0])
    common = len(A[0])
    C = create_matrix(rows, cols)
    for i in range(rows):
        for j in range(cols):
            for k in range(common):
                C[i][j] += A[i][k] * B[k][j]
    return C

# 图像：每个像素是 RGB 三通道 → 就是一个三维矩阵
# 神经网络：权重就是矩阵，前向传播就是矩阵乘法
\end{lstlisting}

\subsection*{练习}

\begin{enumerate}
  \item 用 Python 创建一个 $3 \times 4$ 的矩阵（全部初始化为 0），然后把第二行第三列改为 1。
  \item 有两个矩阵 $A = \begin{bmatrix} 1 & 2 \\ 3 & 4 \end{bmatrix}$ 和 $B = \begin{bmatrix} 5 & 6 \\ 7 & 8 \end{bmatrix}$，手算 $A + B$ 和 $A \times B$。
  \item \textbf{[扩展]} 图像为什么可以用矩阵表示？一张 $100 \times 100$ 的彩色图片，需要多少个数字来存储？
\end{enumerate}

\section*{第十一课：概率初步——不确定性的数学}

\subsection*{随机事件}

\begin{lstlisting}[language=Python]
import random

# 抛硬币
def flip_coin():
    return "正面" if random.random() < 0.5 else "反面"

# 模拟 10000 次
results = {"正面": 0, "反面": 0}
for _ in range(10000):
    results[flip_coin()] += 1
print(f"正面：{results['正面']/10000:.2%}")
print(f"反面：{results['反面']/10000:.2%}")
\end{lstlisting}

\subsection*{期望值}

\begin{lstlisting}[language=Python]
# 掷骰子的期望值
sides = [1, 2, 3, 4, 5, 6]
expected = sum(sides) / len(sides)
print(f"骰子期望值：{expected}")  # 3.5

# 模拟验证
rolls = [random.randint(1, 6) for _ in range(100000)]
print(f"模拟平均：{sum(rolls) / len(rolls):.2f}")
\end{lstlisting}

\subsection*{练习}

\begin{enumerate}
  \item 模拟两个骰子之和的概率分布
  \item 写一个程序验证"生日悖论"：23 人中至少两人同一天生日的概率 > 50\%
  \item \textbf{[扩展]} 用蒙特卡洛方法估算 $\pi$
\end{enumerate}

\section*{第十二课：综合——用数学优化程序}

\subsection*{案例：判断质数}

\begin{lstlisting}[language=Python]
# 方法 1：朴素（O(n)）
def is_prime_naive(n):
    if n < 2:
        return False
    for i in range(2, n):
        if n % i == 0:
            return False
    return True

# 方法 2：优化（O(√n)）
def is_prime_fast(n):
    if n < 2:
        return False
    if n == 2:
        return True
    if n % 2 == 0:
        return False
    # 只需检查到 √n
    i = 3
    while i * i <= n:
        if n % i == 0:
            return False
        i += 2  # 跳过偶数
    return True

# 性能对比
import time
for n in [1000003, 1000033, 1000037]:
    start = time.time()
    r1 = is_prime_naive(n)
    t1 = time.time() - start
    
    start = time.time()
    r2 = is_prime_fast(n)
    t2 = time.time() - start
    
    print(f"{n}: 朴素={t1:.4f}s, 优化={t2:.4f}s, 快了{t1/t2:.0f}倍")
\end{lstlisting}

\subsection*{案例：二分查找的数学原理}

每次查找排除一半数据。$n$ 个元素，需要 $\lceil \log_2 n \rceil$ 次。

\begin{lstlisting}[language=Python]
import math

for n in [10, 100, 1000, 10000, 100000]:
    times = math.ceil(math.log2(n))
    print(f"n={n:6d}: 最多{times:3d}次查找, 而线性查找最多{n}次")
\end{lstlisting}

\subsection*{练习}

\begin{enumerate}
  \item 用数学知识优化以下代码：判断一个数是否是 2 的幂（提示：用位运算，$O(1)$ 时间）。
  \item 对一个长度为 1,000,000 的有序数组做二分查找，最多需要多少次比较？如果改用线性查找呢？
  \item \textbf{[综合]} 你有一个查找系统：先做二分查找（$O(\log n)$），然后做线性扫描（$O(k)$）。对 $n = 10^6$ 和 $k = 100$，总共需要多少次操作？如果把二分查找换成哈希表（$O(1)$ 查找），情况如何？
\end{enumerate}

\subsection*{跨模块关联}

\begin{itemize}
  \item 本模块的\textbf{二进制和位运算}与 \texttt{python-to-cpp/} 中的内存管理和指针密切相关
  \item 本模块的\textbf{图论}与 \texttt{ds-algo/} 中的 BFS/DFS 完全对应
  \item 本模块的\textbf{矩阵和概率}是 \texttt{ai-intro/} 中神经网络的数学基础
  \item 本模块的\textbf{集合论}对应 Python 中的 \texttt{set} 类型（见 \texttt{python/}）
  \item 本模块的\textbf{模运算}在 \texttt{ds-algo/} 的哈希表和循环队列中用到
\end{itemize}

\section*{\textcolor{blue!70!black}{各课练习参考答案}}

\subsection*{第一课：二进制}
\begin{enumerate}
  \item 11010 = 26，27 = 11011。
  \item \texttt{5 \& 3}：0101 \& 0011 = 0001 = 1，与手动计算一致。
  \item 判断 2 的幂：\texttt{n > 0 and (n \& (n-1)) == 0}。因为 2 的幂的二进制只有最高位是 1，减 1 后所有低位变成 1，按位与得 0。
  \item 权限系统：读取=001(1)，写入=010(2)，执行=100(4)。\texttt{permission = read | write = 3}，\texttt{permission \& read != 0} 为 True。
\end{enumerate}

\subsection*{第二课：八进制与十六进制}
\begin{enumerate}
  \item 255 = 0o377 = 0xFF。Python 验证：\texttt{oct(255)} → '0o377'，\texttt{hex(255)} → '0xff'。
  \item \#AABBCC：红色 AA=170，绿色 BB=187，蓝色 CC=204。
  \item 十六进制比八进制更紧凑（每 4 位二进制 vs 3 位），且一个字节（8 位）正好用两个十六进制位表示，非常方便。
\end{enumerate}

\subsection*{第三课：日常生活中的数学概念映射}
\begin{enumerate}
  \item ① \texttt{def f(x): return 2*x + 3} ② \texttt{\{1,2,3\} \& \{3,4,5\}} ③ \texttt{sum(i**2 for i in range(1, 11))}。
  \item 数学上表示为点 $(3, 7)$，Python 中用元组 \texttt{(3, 7)}。
  \item Python 用英文单词而不是数学符号，是为了让代码更接近自然语言、更容易理解和输入。这体现了"可读性优先"的设计原则。
\end{enumerate}

\subsection*{第四课：模运算}
\begin{enumerate}
  \item (3 + 100) \% 7 = 103 \% 7 = 5，所以 100 天后是星期五。
  \item 闰年判断：\texttt{def is\_leap(y): return (y \% 4 == 0 and y \% 100 != 0) or (y \% 400 == 0)}。
  \item 循环队列用 \texttt{(tail + 1) \% size} 实现游标绕回。
\end{enumerate}

\subsection*{第五课：布尔代数}
\begin{enumerate}
  \item XOR 真值表：A=0,B=0→0；A=0,B=1→1；A=1,B=0→1；A=1,B=1→0。
  \item \texttt{not (score >= 60 or grade == "F")} = \texttt{score < 60 and grade != "F"}。
  \item \texttt{A OR B = NOT(NOT A AND NOT B)}，即 \texttt{not (not a and not b)}。
\end{enumerate}

\subsection*{第六课：集合论}
\begin{enumerate}
  \item 综合运用并集、交集、差集：只在一个列表中的元素 = \texttt{(A-B-C) | (B-A-C) | (C-A-B)}。
  \item \texttt{bool(set(str1) \& set(str2))}。
  \item 搜索引擎"与"=交集，"或"=并集，"非"=差集。
\end{enumerate}

\subsection*{第七课：组合数学}
\begin{enumerate}
  \item $10^4 = 10000$ 种。
  \item $C(10, 3) = 120$ 种选法。
  \item \texttt{from itertools import combinations; list(combinations(range(6), 3))} 输出 20 种组合。
\end{enumerate}

\subsection*{第八课：图论基础}
\begin{enumerate}
  \item 5 个顶点的完全图有 $C(5,2) = 10$ 条边。
  \item 无标准答案——用字典表示，如 \texttt{\{"小明": ["小红", "小刚"], ...\}}。
  \item 用 BFS（广度优先搜索）计算无权图的最短路径。
\end{enumerate}

\subsection*{第九课：数列与求和}
\begin{enumerate}
  \item 1 到 1000 所有奇数：共 500 个，和 = $500 \times (1+999)/2 = 250000$。
  \item $2^k \ge n$，$k = \lceil \log_2 n \rceil$。与等比数列 $1, 2, 4, 8, \dots$ 有关，每次翻倍。
  \item 推导：$S = a_1 + a_1q + \dots + a_1q^{n-1}$，两边乘 q 得 $qS = a_1q + \dots + a_1q^n$，相减得 $(q-1)S = a_1(q^n - 1)$，所以 $S = \frac{a_1(q^n - 1)}{q - 1}$。
\end{enumerate}

\subsection*{第十课：矩阵}
\begin{enumerate}
  \item \texttt{m = [[0]*4 for _ in range(3)]; m[1][2] = 1}。
  \item $A+B = \begin{bmatrix}6 & 8 \\ 10 & 12\end{bmatrix}$，$A \times B = \begin{bmatrix}19 & 22 \\ 43 & 50\end{bmatrix}$。
  \item 图像由像素网格组成，每个像素有数值（灰度图一个数，彩色图三个数）。$100 \times 100$ 彩色 = $100 \times 100 \times 3 = 30000$ 个数字。
\end{enumerate}

\subsection*{第十一课：概率初步}
\begin{enumerate}
  \item 两个骰子之和：2-12，概率分布呈三角形，7 的概率最高（1/6）。
  \item 生日悖论验证：随机生成 23 个生日，检查是否有重复，重复实验 10000 次。
  \item 蒙特卡洛估算 $\pi$：在正方形内随机投点，统计落在内切圆中的比例 $\times 4$。
\end{enumerate}

\subsection*{第十二课：综合——用数学优化程序}
\begin{enumerate}
  \item \texttt{n > 0 and (n \& (n-1)) == 0}。时间复杂度 $O(1)$。
  \item $\lceil \log_2 1000000 \rceil = 20$ 次比较。线性查找最多 1,000,000 次。
  \item 二分+线性：$20 + 100 = 120$ 次操作。哈希+线性：$1 + 100 = 101$ 次操作。哈希表略快但占用更多内存。
\end{enumerate}

\end{document}
